\chapter{conditional distribution and conditional expectation}
\section{random vector and joint distributions}
\subsection{random vector}
\begin{definition}{random vector}
对于 prob space $(\Omega, \mathcal{F}, P)$, 
一个 function $\mathbf{X}: \Omega \to \mathbb{R}^n$ 如果是一个
$(\mathcal{F}, \mathcal{B}(\mathbb{R}^n))$-measurable function, 则
称它为一个 $n$-dimensional random variable, 或者 $n$-dimensional random vector.

通常我们将 random vector 写成分量形式
$$\mathbf{X}(\omega) = (X_1(\omega), X_2(\omega), \dots, X_n(\omega))^T$$
\end{definition}
random vector 相当于在一个 prob space 上, 考虑多个重新分配 mass 的方法, 并把它们并列起来.
\begin{proposition}{由 \(n\) 个 random variable 构成的 vector 是一个 random vector}
    令 \(X_1, X_2, \cdots, X_n\) 
    是定义在\textbf{同一个 prob space} \((\Omega, \mathcal{F}, \bP)\) 
    上的 \(n\) 个 random variables. 则函数 \(\mathbf{X}: \Omega \to \bR^n\) defined by 
    \[
    \omega \mapsto (X_1(\omega), X_2(\omega), \cdots, X_n(\omega))^T
    \]
    是一个 random vector.
\end{proposition}
\begin{proof}
    我们在 measure theory 中证明过: 对于任意的 finite seq of Borel measureable functions 
    \( (  {f_i}: \Omega \to \mathbb{R} )_{i=1}^k\), 其各作为一个维度组成的函数
    \(f = (f_1,\cdots, f_k)\) 也是一个 Borel measurable function (from \(\Omega\) 到 \(\mathbb{R}^k\)).
\end{proof}
\begin{proposition}{一个 random vector 的每个分量都是一个 random variable}
    令 \(\mathbf{X} = (X_1, X_2, \cdots, X_n)^T\) 是一个 random vector. 
    则对于任意 \(i\), \(X_i\) 都是一个 random variable.
\end{proposition}
\begin{proof}
    对于 $\mathbf{X}: \Omega \to \mathbb{R}^n$ 是 $(\mathcal{F}, \mathcal{B}(\mathbb{R}^n))$-measurable,
    我们可以将第 $i$ 个分量 $X_i$ 看作是 composition of two maps: 
    $$X_i = \pi_i \circ \mathbf{X}$$
    其中 $\pi_i: \mathbb{R}^n \to \mathbb{R}$ projection map  $\pi_i(x_1, \dots, x_n) = x_i$.\\
    由于 projection map $\pi_i$ 是一个 Borel measurable function, 我们可以得出: \(X_i = \pi_i(\mathbf{X})\) 也是
一个 Borel measurable function, 也就是一个 random variable.
\end{proof}

\begin{remark}
上两条 propositions 说明了一个事情: random vectors 和其分量 random variables 完全相互决定.
\begin{itemize}
    \item\textbf{ random variables 组合起来一定是一个 random vector.}
    \item \textbf{random vector 拆开每个分量一定是 random variable.}
\end{itemize}
因而, 研究一个 random vector 也就相当于研究分量 random variables 之间的关系. 

这一章节我们会从 random vector 的性质 
(比如 \textbf{joint distribution 和 marginal distribution, conditional distribution}) 
出发, 研究分量 random variables 之间的关系.

(但是注意: 前提是这些 random variables 必须定义在同一个 prob space 上. 也就是, 它们的 base 
probability measure 一定相同. 如果我们考虑不同 prob space 上的 random variables 的话, 
那它们没有一个公共的 base probability measure, 很难说它们之间有什么关系. 
这种情况下, 我们需要 coupling: 构造一个新的 prob space (说白了就是 product space),
让这两个 random variables 都有一个新 version. 这一问题我们之后再讨论.)
\end{remark}

\subsection{joint distribution}
\begin{definition}{joint distribution and joint cdf}
    令 \(X_1,\cdots, X_n\) 为 RV from the same prob space. 
    即 \(\mathbf{X}:= (X_1,\cdots, X_n)\) 
    是一个 random vector.

    我们称 \(\bP^{\mathbf{X}}: \bR^n \to \bR\) defined by 
    $$\bP^{\mathbf{X}}(B) = \bP(\mathbf{X}^{-1}(B)), \quad \forall B \in \mathcal{B}(\mathbb{R}^n)$$
    为 \(X_1, \cdots, X_n\) 的 \textbf{joint distribution}.
    
    而我们称 \(F_{\mathbf{X}} = F_{X_1,
    \cdots, X_n}: 
    \bR^n \to \bR\) defined by   
     \[
    F_{\mathbf{X}}(x_1,\cdots,x_n) = \bP(X_1 \leq x_1, \cdots, X_n \leq x_n)
    \]
     (即 \(\bP^{\mathbf{X}}\) restricted to the set of rectangles
     \(\{(-\infty, x_1] \times \cdots \times (-\infty, x_n]: (x_1,\cdots,x_n) \in \bR^n\}\))
    为它们的 \textbf{joint distribution function (或称 joint cdf)}.
\end{definition}
joint distribution 的定义已经包括了如何从多个 random variables 的 distributions
得到一个 joint distribution.
而, 我们也可以从一个 joint distribution 的 limit behavior 
得到每个 random variable 分量的 distributions, 称之为 marginal distribution:
\begin{definition}{marginal distribution}
    令 \(\mathbf{X} = (X_1, \cdots, X_n)^T\) 是一个 random vector. 
    则对于任意 \(i\), \(X_i\) 的 distribution \(\bP^{X_i}\) 
    被称为 \(\mathbf{X}\) 的第 \(i\) 个分量的 \textbf{marginal distribution}.
\end{definition}
我们以 \(\bR^2\) 为例. 得出的结论可以推广到 \(\bR^n\).
\begin{proposition}{通过 joint distribution 的 limit behavior 得到 marginal distribution}
    令 \(\mathbf{X} = (X_1, X_2)^T\) 是一个 random vector. 
    则对于任意 \(x\in \bR\), 有
    \[\bP(X_1 \leq x) = \lim_{y\to \infty} F_{\mathbf{X}}(x,y)\]
    \(X_2\) 的 marginal distribution 也可以通过同样的方法得到.

    此外, joint distribution 有其他明显的 limit behaviors:
    \begin{itemize}
        \item \[ \lim_{x\to -\infty}F(x,y) = \lim_{y\to -\infty}F(x,y) = 0 \]
        \item \(F_{\mathbf{X}}\) 对于每个维度都是 increasing 且 right-continuous 的.
        \item \[ \bP(X\leq x, Y \leq y) = \lim_{z\uparrow y} F(x,z) = \lim_{z\uparrow x} F(z,y) = \lim_{z_1\uparrow x, z_2 \uparrow y}F(z_1, z_2) \]
        \item \[ \bP(x_1 \leq X \leq x_2, y_1 \leq Y \leq y_2) = F(x_2,y_2) - F(x_1,y_2) = F(x_2,y_1)+F(x_1,y_1) \]
    \end{itemize}
\end{proposition}
\begin{remark}
    最后一条: 右边相当于 
    \[
    \bP(x_1 < X < x_2, Y < y_2) - \bP(x_1 < X < x_2, Y < y_1) 
    \]
    因而成立.
\end{remark}

discrete random vector 很容易处理. 我们可以直接定义 joint pmf.
而 continuous random vector 需要展开讨论. 接下来我们讲单独讨论
continuous random vector 的 joint distribution.

\subsection{condinuous joint cdf 与 joint pdf}
recall:
 一个随机变量 \(X\) 被称为一个 continuous random variable, 如果它的
cdf 是 absolutely continuous 的. 这个条件也等价于, 存在一个函数 \(f_X: \bR \to [0,\infty)\) 
使得 \[
F_X(x) = \int_{-\infty}^x f_X(t) \, dt
\]

这个定义可以 generalize 到 random vector 上.
\begin{definition}{continuous random vector 和 continuous joint cdf}
    对于 random vector  \(\mathbf{X} = (X_1,\cdots, X_n)^T\) 
    如果 \(\bP^{\mathbf{X}} \ll \lambda^n\),
    即存在一个函数 \(f_{\mathbf{X}}: \bR^n \to [0,\infty)\) 使得对于任意 Borel set \(B \subseteq \bR^n\), 有
\[\bP^{\mathbf{X}}(B) = \int_B f_{\mathbf{X}}(x_1, \cdots, x_n) \, d\lambda^n(x_1, \cdots, x_n)\]
    则称 \(\mathbf{X}\) 是一个 \textbf{continuous random vector}.

    注意: 由于是在 \(\bR^n\) 上, 这等价于存在一个函数 \(f_{\mathbf{X}}: \bR^n \to [0,\infty)\) 使得对于任意 \(x_1, \cdots, x_n\), 有
     \[F_{\mathbf{X}}(x_1, \cdots, x_n) = \int_{-\infty}^{x_1} \cdots \int_{-\infty}^{x_n} f_{\mathbf{X}}(t_1, \cdots, t_n) \, dt_n \cdots dt_1\]
\end{definition}
\begin{remark}
    这个函数 \(f_{\mathbf{X}}\) 
    就是 continuous random vector \(\mathbf{X}\) 
    的 \textbf{joint probability density function (joint pdf).}
    它是 joint distribution measure \(\bP^{\mathbf{X}}\) 
    对于 Lebesgue measure \(\lambda^n\) 的 Radon-Nikodym derivative: \[\frac{d\bP^{\mathbf{X}}}{d\lambda^n} = f_{\mathbf{X}}\]
\end{remark}
\begin{remark}
    我们已经知道, 只要存在导数 \(f_{\mathbf{X}}\) 可以对于任意 \(\mathbf{x}\in \bR^n\) 
    还原出 joint cdf \(F_{\mathbf{X}}\) 的值, 那么 \(f_{\mathbf{X}}\) 就是 joint pdf, 
    这个 random vector 就是一个 continuous random vector.

    那么这个还原积分的计算可以任意换序吗? 

    显然是可以的. 因为 recall \textbf{Fubini's theorem}: 只要 \(f\) 绝对可积, 即 
    \(\int_{\bR^n} |f| \, d\lambda^n < \infty\), 
    那么对于任意的 permutation \(\sigma\) of \(\{1,2,\cdots,n\}\), 我们都可以将积分的顺序换成
     \[\int_{-\infty}^{x_{\sigma(1)}} \cdots \int_{-\infty}^{x_{\sigma(n)}} 
     f(t_1, \cdots, t_n) \, dt_{\sigma(n)} \cdots dt_{\sigma(1)}\]
    而我们知道, \(f\) 一定是非负的, 而且积分一定是 1, 因而可以任意换序积分. 
    
    这就很舒服. 这也给出了 margin distribution 的计算方法: 
    只要对 joint pdf \(f_{\mathbf{X}}\) 在其他维度上积分掉就行了.
\end{remark}


\begin{proposition}{joint pdf 和 joint cdf 的性质}
令 \(\mathbf{X} = (X_1, \cdots, X_n)^T\) 是一个 continuous random vector, 
则它的 joint pdf \(f_{\mathbf{X}}\) 和 joint cdf \(F_{\mathbf{X}}\) 有以下性质: 
    \begin{itemize}
        \item \(f_{\mathbf{X}}\geq 0\) a.e. 并且 \(\int_{\bR^n} f_{\mathbf{X}} \, d\lambda^n = 1\).
\item (如果一个集合 \(A\) 有一个零测维度, 那么 \(\bP^{\mathbf{X}}(A) = 0\))
\[\bP(X_1 = x_1
, x_2 \leq X_2 \leq x'_2  \cdots,  x_n \leq  X_n \leq x'_n) = 0\]

\item 每个 \(X_i\) 的 marginal distribution 也是 (absolutely) continuous 的,
并且 \[
f_{X_i}(x) = \int_{\bR^{n-1}} f_{\mathbf{X}}(t_1, \cdots, t_{i-1}, x, t_{i+1}, \cdots, t_n) \, dt_1 \cdots dt_{i-1} dt_{i+1} \cdots dt_n
\]
例如, \[
f_{X}(x) = \int_{-\infty}^\infty f_{\mathbf{X}}(x,y) \, dy
\]
\item 对每个 \(x_1, \cdots, x_n \in \bR\), 都可以通过偏导数从 joint cdf 得到 
joint pdf (这个偏导数一定 (a.e.) 存在):
\[
f_{\bf{X}}(x_1,\cdots,x_n) = \frac{\partial^n}{\partial x_1 \cdots \partial x_n} F_{\mathbf{X}}(x_1, \cdots, x_n)
\]
并且 by Fubini's theorem 可以 (a.e.) 任意换序:
$$\frac{\partial^n}{\partial x_{\sigma(1)} \dots \partial x_{\sigma(n)}} F_{\mathbf{X}} \stackrel{a.e.}{=} f_{\mathbf{X}}$$
\end{itemize}
\end{proposition}
\begin{proof}

\begin{itemize}
    \item 由于 \(f_{\mathbf{X}}\) 是 \(\bP^{\mathbf{X}}\) 对于 \(\lambda^n\) 的 Radon-Nikodym derivative, 
    因为任何 Borel 集 $B$ 都有 $P^{\mathbf{X}}(B) \ge 0$,
     假设存在一个集合 $A \in \mathcal{B}(\mathbb{R}^n)$ 
     使得在其上 $f_{\mathbf{X}} < 0$ 且 $\lambda^n(A) > 0$, 
    那么根据积分定义$$\bP^{\mathbf{X}}(A) = \int_A f_{\mathbf{X}} \, d\lambda^n < 0$$
    因而反证得 \(f_{\mathbf{X}} \geq 0\) a.e. 

    并且 \[\int_{\bR^n} f_{\mathbf{X}} \, d\lambda^n = \bP^{\mathbf{X}}(\bR^n) = 1\]

\item natural.
\item 即 marginal distribution 的定义在 continuous case 中的展开
\item by Fubini's theorem, 以及 joint cdf 的定义.
\end{itemize}
    \end{proof}

\begin{example}
考虑 
\[
f_{X, Y}(x, y)= \begin{cases}c e^{-x} e^{-2 y}, & \text { if } x, y>0 \\ 0, & \text { otherwise }\end{cases}
\]
Find \(c\) 使得这是一个合法的 joint pdf, 
并且计算 \(X\) 和 \(Y\) 的 marginal pdfs, 以及 \(\bP(X>1, Y<1)\).
\begin{solution}
我们需要
$$c \left( \int_{0}^{\infty} e^{-x} dx \right) \left( \int_{0}^{\infty} e^{-2 y} dy \right) = 1$$
evaluate 这两个积分:
$$\int_{0}^{\infty} e^{-x} dx = \left[ -e^{-x} \right]_{0}^{\infty} = 1,$$
$$\int_{0}^{\infty} e^{-2 y} dy = \left[ -\frac{1}{2} e^{-2 y} \right]_{0}^{\infty} = \frac{1}{2}$$
因为 \(c\) 必须为 2.

计算 \(X\) 的 marginal pdf:
$$f_X(x) = \int_{0}^{\infty} 2 e^{-x} e^{-2 y} dy = 2 e^{-x} \left[ -\frac{1}{2} e^{-2 y} \right]_{0}^{\infty} = e^{-x}$$
然后计算 \(Y\) 的 marginal pdf:
$$f_Y(y) = \int_{0}^{\infty} 2 e^{-x} e^{-2 y} dx = 2 e^{-2 y} \left[ -e^{-x} \right]_{0}^{\infty} = 2 e^{-2 y}$$
最后计算 \(\bP(X>1, Y<1)\):
$$\bP(X > 1, Y < 1) = \int_{1}^{\infty} \int_{0}^{1} 2 e^{-x} e^{-2 y} dy dx = \left( \int_{1}^{\infty} e^{-x} dx \right) \left( \int_{0}^{1} 2 e^{-2 y} dy \right)$$
这两个定积分分别为:
$$\int_{1}^{\infty} e^{-x} dx = \left[ -e^{-x} \right]_{1}^{\infty} = e^{-1},$$
$$\int_{0}^{1} 2 e^{-2 y} dy = \left[ -e^{-2 y} \right]_{0}^{1} = 1 - e^{-2}$$ 
因而$$\bP(X > 1, Y < 1) = e^{-1} (1 - e^{-2}) = e^{-1} - e^{-3}$$
\end{solution}
\end{example}


\section{independence and correlation}
\subsection{independence}
\begin{definition}{independence of random variables}
    l两个 random variables \(X, Y: \Omega \to \bR\) 被称为 independent 的, 
    如果对于任意的 Borel sets \(A, B \subseteq \bR\), 都有
    \[
    \bP(X \in A, Y \in B) = \bP(X \in A) \cdot \bP(Y \in B)
    \]
    即: \textbf{它们的 joint distribution 是它们 marginal distributions 的 product.}
    注意, 这个定义等价于 \[
    F_{X,Y} = F_X F_Y
    \]
\end{definition}

\begin{proposition}{discrete and continuous independence: joint density is product of marginal densities} \label{discrete and continuous case of independence: joint density is product of marginal densities}
    令 \(X, Y\) 是两个 random variables. 
    \begin{itemize}
        \item 如果 \(X, Y\) 是 discrete random variables, 
        则它们 independent 的条件等价于对于任意 \(x, y\), 都有
        \[
        p_{X,Y} = p_X \cdot p_Y
        \]
        \item 如果 \(X, Y\) 是 continuous random variables, 则它们 independent 的条件等价于对于任意 \(x, y\), 都有
        \[
        f_{X,Y} = f_X \cdot f_Y
        \]
    \end{itemize}
\end{proposition}
\begin{proof}
    discrete case: 显然可得.
    continuous case: 
    \(F_{X,Y} = F_X F_Y  \implies f_{X,Y} = f_X \cdot f_Y\): 取偏导数即可.
    \(f_{X,Y} = f_X \cdot f_Y \implies F_{X,Y} = F_X F_Y\): 积分即可.
\end{proof}

\subsection{covariance and correlation}


\section{conditional distribution}
